%==============================================================================
\chapter{Methodology}
\label{ch:methodology}
%==============================================================================

This chapter presents the methodology of the work: the requirements that
define the problem precisely (Section~\ref{sec:requirements}), the three
locked design decisions that shape the whole system
(Section~\ref{sec:decisions}), the agentic generation method with its
deterministic validation gate (Section~\ref{sec:genmethod}), the
synchronization model (Section~\ref{sec:syncmethod}), the caching and
idempotency strategy (Section~\ref{sec:cachemethod}), and the phased
development and testing process followed (Section~\ref{sec:sdlc}).

%------------------------------------------------------------------------------
\section{Requirements Analysis}
\label{sec:requirements}
%------------------------------------------------------------------------------

\subsection{Target Users and User Stories}

The system targets three user groups: students revising a concept,
self-learners exploring an unfamiliar topic, and educators generating quick
visual aids. The core user stories are listed in Table~\ref{tab:stories}.

\begin{table}[htbp]
\caption{Core user stories}
\label{tab:stories}
\centering
\small
\renewcommand{\arraystretch}{1.3}
\begin{tabular}{@{}lp{0.14\textwidth}p{0.34\textwidth}p{0.30\textwidth}@{}}
\toprule
\textbf{Id} & \textbf{As a\ldots} & \textbf{I want to\ldots} & \textbf{So that\ldots}\\
\midrule
US-1 & learner & type a topic and get a narrated diagram & I understand it visually and aurally\\
US-2 & learner & see diagram parts highlight as they are explained & I know exactly what is being described\\
US-3 & learner & pause, resume, and replay any segment & I can learn at my own pace\\
US-4 & learner & find all my past topics in a history list & I can revisit lessons\\
US-5 & learner & replay a saved lesson with no internet & I can study offline\\
US-6 & learner & retry generation if a result is poor & I get a usable lesson\\
\bottomrule
\end{tabular}
\end{table}

\subsection{Functional Requirements}

\begin{description}
  \item[FR-1] Accept a free-text topic from the user.
  \item[FR-2] Generate an SVG diagram with stable, semantic element
        identifiers.
  \item[FR-3] Generate an ordered narration script split into segments, each
        referencing the SVG element identifiers it describes.
  \item[FR-4] Produce one TTS audio clip per segment with a measured
        duration.
  \item[FR-5] Assemble a self-contained lesson manifest (SVG + segments +
        audio references + durations).
  \item[FR-6] Serve cached lessons for previously requested topics.
  \item[FR-7] On device: download and persist the full lesson (manifest +
        assets) locally.
  \item[FR-8] On device: play the lesson, highlighting each segment's SVG
        element identifiers while its audio plays.
  \item[FR-9] On device: provide play, pause, resume, seek-by-segment, and
        replay controls.
  \item[FR-10] On device: maintain a searchable history of saved lessons.
  \item[FR-11] On device: replay any saved lesson fully offline.
  \item[FR-12] Surface clear errors and allow regeneration on failure.
\end{description}

\subsection{Non-Functional Requirements}

\begin{table}[htbp]
\caption{Non-functional requirements}
\label{tab:nfr}
\centering
\small
\renewcommand{\arraystretch}{1.3}
\begin{tabular}{@{}llp{0.62\textwidth}@{}}
\toprule
\textbf{Id} & \textbf{Category} & \textbf{Requirement}\\
\midrule
NFR-1 & Performance & Cache hit returns a manifest in under 500\,ms; cold
generation target under 30\,s\\
NFR-2 & Offline & Saved lessons play with zero network calls\\
NFR-3 & Reliability & Generation validates SVG well-formedness and identifier
consistency before caching\\
NFR-4 & Cost & Identical topics are generated once and shared via cache\\
NFR-5 & Portability & LLM and TTS providers are swappable behind adapters\\
NFR-6 & Maintainability & Clear layering, documented design patterns,
unit-testable core\\
NFR-7 & Privacy & No accounts; personal history never leaves the device\\
NFR-8 & Accessibility & Captions shown; adjustable playback; readable
contrast\\
\bottomrule
\end{tabular}
\end{table}

\subsection{Scope Exclusions (Version 1)}

User accounts and cross-device cloud synchronization, editing or annotating
generated diagrams, a multi-language user interface (lesson language is a
parameter, but the UI is single-language), social and sharing features, and
real-time word-level highlighting (segment-level synchronization is used, as
motivated in Section~\ref{sec:syncmethod}) are explicitly out of scope.

%------------------------------------------------------------------------------
\section{Key Design Decisions}
\label{sec:decisions}
%------------------------------------------------------------------------------

Three decisions, taken early and documented before implementation, shape
everything else in the system.

\textbf{Decision 1: Synchronization by pre-rendered timed segments.} Rather
than align voice to highlight at play time, the server splits the narration
into segments, renders one audio clip per segment, and \emph{measures} each
clip's duration. The device simply plays the clips in order and highlights
the parts named by the current clip. This makes synchronization deterministic
and removes any need for the network during playback (NFR-2), at the cost of
segment-level rather than word-level pointer granularity.

\textbf{Decision 2: Stateless backend with a shared cache.} There are no user
accounts on the server (NFR-7). Lessons are cached under a normalized topic
key, so a popular topic is generated once and then served to everyone
(NFR-1, NFR-4). Statelessness keeps the server trivially scalable and simple
to reason about; all personal state (history) lives on the device.

\textbf{Decision 3: Provider adapters.} The language model and the speech
engine sit behind narrow interfaces; the pipeline talks to those interfaces
and never to a vendor SDK directly (NFR-5). Choosing a model --- or replacing
the vendor entirely --- is a configuration change. This also acknowledges an
operational reality: hosted model identifiers change over time, so the model
name is treated as configuration, verified against the live API by a small
helper tool, and pinned in settings rather than in code.

%------------------------------------------------------------------------------
\section{The Agentic Generation Method}
\label{sec:genmethod}
%------------------------------------------------------------------------------

\subsection{Why a Graph and Not a Single Prompt}

A single mega-prompt (``produce a diagram, narration, and timing for topic
X'') couples all failure modes together: if any part of the output is
malformed, the whole response must be regenerated, and there is no precise
error signal to steer the retry. Decomposing generation into a
\emph{state graph} of small steps gives three properties the product needs:

\begin{enumerate}
  \item each artifact (plan, drawing, narration) can be \emph{validated and
        repaired in isolation}, with concrete error messages;
  \item each node is independently testable against mock providers; and
  \item deterministic steps (validation, TTS rendering, assembly) are
        separated from probabilistic ones, so reliability mechanisms attach
        exactly where the risk is.
\end{enumerate}

The pipeline is implemented as a LangGraph \texttt{StateGraph}: a single
typed state object flows through nodes connected by conditional edges. The
graph is small, static, and every loop carries a configured budget --- this
is deliberate rejection of open-ended agent autonomy in favour of an
inspectable, bounded process.

\subsection{Pipeline Stages}

Figure~\ref{fig:pipeline} shows the graph. Its nodes are:

\begin{figure}[htbp]
\centering
\begin{tikzpicture}[
  font=\small,
  node distance=5mm,
  n/.style={draw, rounded corners, align=center, minimum width=62mm,
            minimum height=8mm, fill=blue!4},
  d/.style={draw, diamond, aspect=2.2, align=center, inner sep=1pt,
            fill=yellow!12},
  >={Stealth}
]
\node[n] (plan) {planner: outline concepts};
\node[n, below=of plan] (svg) {svg\_generator: draw SVG with ids};
\node[d, below=of svg] (crit) {critique\\good?};
\node[n, below=of crit] (narr) {narration: segments $\rightarrow$ ids};
\node[d, below=of narr] (val) {validate\\ids exist?};
\node[n, below=of val] (tts) {tts\_renderer: audio + duration};
\node[n, below=of tts] (asm) {assembler: build + cache};

\draw[->] (plan) -- (svg);
\draw[->] (svg) -- (crit);
\draw[->] (crit) -- node[right,font=\footnotesize]{yes} (narr);
\draw[->] (narr) -- (val);
\draw[->] (val) -- node[right,font=\footnotesize]{yes} (tts);
\draw[->] (tts) -- (asm);

\draw[->] (crit.west) to[out=180,in=180,looseness=1.5]
  node[left,font=\footnotesize,align=center]{no:\\refine\\(bounded, $K$)} (svg.west);
\draw[->] (val.east) to[out=0,in=0,looseness=1.5]
  node[right,font=\footnotesize,align=center]{no:\\repair\\(bounded, $R$)} (narr.east);
\end{tikzpicture}
\caption{The generation graph. A bounded critique-and-refine loop improves
the drawing before narration is written; a bounded repair loop fixes
narration that fails the deterministic validator. A lesson reaches the
assembler only after it passes validation.}
\label{fig:pipeline}
\end{figure}

\begin{description}
  \item[Planner.] The model outlines the concepts and a visual plan for the
        topic. Planning once, up front, keeps the drawing and the narration
        aligned to the same concept set.
  \item[SVG generator.] Given the plan, the model draws the SVG and assigns a
        semantic \texttt{id} to every meaningful part (for example
        \texttt{sun}, \texttt{leaf}, \texttt{chloroplast}). The prompt
        instructs the model to return \emph{only} SVG.
  \item[Critique and refine.] A critique step scores the drawing and lists
        concrete problems. If the score is below a threshold and budget
        remains, a refine step redraws with that feedback and the critique
        runs again, up to $K$ passes. If the budget runs out, the best
        drawing so far is accepted. This is a focused application of the
        self-refine idea~\cite{madaan2023selfrefine}, applied only to the
        artifact that benefits most.
  \item[Narration generator.] Given the \emph{finished} drawing, the model
        writes ordered narration as schema-constrained JSON: each segment has
        text and a list of element identifiers. Writing narration against the
        real drawing steers the model toward referencing parts that exist ---
        but does not guarantee it, which is why the next node exists.
  \item[Validator.] This node uses no model (Section~\ref{sec:gate}).
  \item[Repair.] If validation fails and retries remain (budget $R$), the
        repair step receives the \emph{exact} list of offending identifiers
        and parse errors and is asked for the minimal correction; validation
        then runs again. If the budget is exhausted, the job fails cleanly
        --- the system refuses service rather than caching a broken lesson.
  \item[TTS renderer.] For each segment, the speech provider synthesizes
        audio; the clip is stored as an asset and its measured duration
        recorded. Segments are independent, so this step is parallelizable.
  \item[Assembler.] A builder assembles the lesson incrementally --- title,
        SVG asset, and each segment with text, identifiers, audio asset, and
        duration --- and the finished lesson is cached.
\end{description}

\subsection{The Deterministic Validation Gate}
\label{sec:gate}

The validator is the reason the product's central feature is trustworthy, so
its checks are stated precisely. Given the SVG string and the narration
segments, it verifies, using only an XML parser and simple logic:

\begin{enumerate}
  \item the SVG parses as well-formed XML and declares a \texttt{viewBox};
  \item every identifier referenced by every segment's
        \texttt{svg\_element\_ids} exists as an \texttt{id} attribute in the
        SVG;
  \item narration is non-empty, segment texts are non-empty, and segments
        are ordered and contiguous from index~0; and
  \item (as a soft, non-blocking signal) major SVG elements never referenced
        by any segment are flagged for coverage.
\end{enumerate}

The result is a pass-or-fail verdict with a machine-readable error list.
Because the validator contains no model, it cannot hallucinate, costs
nothing, and always terminates. It is the \emph{hard gate}: the invariant
enforced system-wide is that \emph{no lesson whose narration references a
non-existent element is ever cached or served}. LLM-based judges were
deliberately not used for this check --- a probabilistic judge in front of a
correctness property merely moves the hallucination problem, whereas an XML
parser settles it.

\subsection{Prompting Strategy}

Four practices keep the probabilistic steps steerable: (i) the \emph{stable
ids contract} embedded in the SVG and narration prompts; (ii)
\emph{structured output} everywhere a machine consumes the response, with
Pydantic validation at the boundary; (iii) \emph{repair prompts} that contain
the validator's concrete error list rather than a vague ``fix this''; and
(iv) \emph{low temperature} for validation-sensitive steps, with bounded
retries ($K$ and $R$, typically 2) to cap cost and latency.

%------------------------------------------------------------------------------
\section{The Synchronization Model}
\label{sec:syncmethod}
%------------------------------------------------------------------------------

Let a lesson have segments $s_0, s_1, \ldots, s_{n-1}$, where each $s_i$
carries text, a set of element identifiers $E_i$, and an audio clip of
measured duration $d_i$ milliseconds. The total lesson duration is
$D=\sum_{i} d_i$. Playback is a simple loop driven by clip boundaries: when
clip $i$ starts, highlight exactly the elements in $E_i$; when it ends, clear
the highlight and advance to $i+1$. Pause, resume, and seek-by-segment fall
out naturally, because the only synchronization state is the index of the
current clip --- a quantity the audio player maintains anyway.

Because each $d_i$ is computed on the server from the synthesized audio's
byte count and sample rate (not estimated), the on-device timing cannot
drift: the highlight changes precisely when the audio item changes. No
network, no timestamps, and no alignment service are involved, so the same
mechanism works identically online and offline. The trade-off, accepted
knowingly, is granularity: a group of parts is lit for a whole sentence
rather than each part at its spoken word.

%------------------------------------------------------------------------------
\section{Caching, Idempotency, and the Job Model}
\label{sec:cachemethod}
%------------------------------------------------------------------------------

Generation is expensive (seconds, several model calls) while a cache hit is
cheap, so the API is asynchronous. A topic is normalized into a
\emph{topic key} (lower-cased, trimmed, whitespace-collapsed, language
suffixed --- for example \texttt{photosynthesis|en}), which serves three
roles at once:

\begin{itemize}
  \item \textbf{cache key} --- a request for a cached topic returns the
        manifest immediately;
  \item \textbf{idempotency key} --- repeated \texttt{POST} requests for the
        same topic return the same job or the same cached lesson; and
  \item \textbf{coalescing key} --- concurrent requests for the same topic
        attach to one in-flight job instead of spawning duplicate
        generations, so the peak cost of a suddenly popular topic is one
        generation.
\end{itemize}

A generation job progresses through
$\texttt{queued} \to \texttt{running} \to \texttt{succeeded} \mid
\texttt{failed}$, exposing progress and the current pipeline stage; the
client polls with backoff until a terminal state
(Section~\ref{sec:apidesign}).

%------------------------------------------------------------------------------
\section{Development Methodology and Testing Strategy}
\label{sec:sdlc}
%------------------------------------------------------------------------------

\subsection{Phased, Design-First Process}

The project followed a documented, phased process with an explicit exit
criterion per phase (Table~\ref{tab:phases}) --- design first (a full set of
UML artifacts: component, class, ER, sequence, and state diagrams), then
vertical implementation slices, each leaving the system demonstrably working.
Two process rules were applied throughout. First, \emph{providers are
mockable from day one}: every backend phase runs against deterministic mock
LLM/TTS implementations, so development and tests never depend on live AI
keys. Second, \emph{definition of done} per phase includes code, tests, and
updated documentation and diagrams.

\begin{table}[htbp]
\caption{Development phases and exit criteria}
\label{tab:phases}
\centering
\small
\renewcommand{\arraystretch}{1.3}
\begin{tabular}{@{}clp{0.55\textwidth}@{}}
\toprule
\textbf{Phase} & \textbf{Name} & \textbf{Exit criterion}\\
\midrule
0 & Vision \& requirements & Requirements agreed\\
1 & Architecture \& design & UML design set reviewed and approved\\
2 & Backend skeleton & \texttt{POST /lessons} returns a mock lesson end to end\\
3 & AI pipeline & Graph produces a valid lesson from mocks; validator rejects
bad identifiers\\
4 & Real generation & Real topic $\to$ cached lesson with audio; cache-hit
path works\\
5 & Android skeleton & App generates a lesson and shows its manifest\\
6 & Player + highlight & Narration plays with parts highlighting in sync\\
7 & Offline \& history & Saved lesson replays with airplane mode on\\
8 & Hardening \& CI/CD & Green CI; hermetic tests; containerized backend\\
9 & Release preparation & Installable build + deployable backend\\
\bottomrule
\end{tabular}
\end{table}

\subsection{Testing Strategy}

Testing is organized by level on both sides of the system
(Table~\ref{tab:testing}). The distinctive property of the backend suite is
that it is \emph{hermetic and key-free}: a session-level fixture disables
environment-file loading and strips all project environment variables, so the
full suite runs with zero live API calls and no secrets --- in any clone and
in CI. The Gemini adapters themselves are tested against a faked SDK client,
which exercises the adapter logic (structured output handling, PCM-to-WAV
wrapping, duration computation, retry/backoff) without the network. Risk
analysis and mitigations tracked during development are summarized in
Table~\ref{tab:risks}.

\begin{table}[htbp]
\caption{Testing levels}
\label{tab:testing}
\centering
\small
\renewcommand{\arraystretch}{1.3}
\begin{tabular}{@{}lp{0.40\textwidth}p{0.38\textwidth}@{}}
\toprule
\textbf{Level} & \textbf{Backend} & \textbf{Android}\\
\midrule
Unit & Graph nodes, validator, builder, providers (mocked) & ViewModels,
use cases, mappers\\
Contract & API request/response schema tests & Repository $\leftrightarrow$
API contract\\
Integration & Graph end to end with mock providers; cache and asset store &
Room DAO tests\\
End to end & Golden-topic generation smoke against live models & Compose UI
and offline-playback instrumentation\\
\bottomrule
\end{tabular}
\end{table}

\begin{table}[htbp]
\caption{Principal risks and mitigations}
\label{tab:risks}
\centering
\small
\renewcommand{\arraystretch}{1.3}
\begin{tabular}{@{}p{0.30\textwidth}p{0.62\textwidth}@{}}
\toprule
\textbf{Risk} & \textbf{Mitigation}\\
\midrule
SVG quality/consistency from the LLM & Plan $\to$ generate $\to$ critique
$\to$ refine loop; low temperature; golden-topic checks\\
Identifier--narration mismatch & Hard validator gate: never cache if any
referenced identifier is missing\\
Hosted model identifier drift & Model name is configuration; helper verifies
against the live API; adapter isolates the SDK\\
TTS synchronization drift & Segment-level clips with byte-measured durations
make synchronization deterministic\\
Cold-start latency & Asynchronous jobs with progress UI; shared cache;
per-segment TTS parallelizable\\
Cost per generation & Topic-keyed cache; request coalescing; bounded loops\\
Injected script in generated SVG & Treat SVG as untrusted; WebView restricted
to the app's own highlight script\\
Renderer choice (WebView vs.\ native) & Highlight engine behind an interface;
decided by a spike in Phase~6\\
\bottomrule
\end{tabular}
\end{table}
